\makeatletter
\@ifundefined{keepfig}{\newenvironment{keepfig}{\par\edef\keepfig@pi{\the\parindent}\noindent
  \begin{minipage}{\linewidth}\parindent=\keepfig@pi\relax}{\end{minipage}\par}}{}
\makeatother
\section{Группа Галуа многочлена}

\begin{definition}[3.27a]
Пусть $f\in F[X]$~--- сепарабельный многочлен, $F_f$~--- его поле разложения. Тогда $F_f/F$~--- расширение Галуа
(\rf{3.10}), и группа $G_f=\Gal(F_f/F)$ называется \emph{группой Галуа многочлена $f$}.
\end{definition}

Поле разложения определено с точностью до $F$-изоморфизма (\rf{2.13}), и изоморфизм $\varphi\colon F_f\to F_f'$ полей
разложения переводит группу $\Gal(F_f/F)$ в $\Gal(F_f'/F)$ по правилу $\sigma\mapsto\varphi\sigma\varphi^{-1}$; так что
$G_f$ определена с точностью до изоморфизма. Если $F$ совершенно (например, $\operatorname{char}F=0$), то группу Галуа
можно определить для любого $f\ne0$: поле разложения $f$ совпадает с полем разложения произведения различных
унитарных неприводимых множителей $f$, а этот многочлен сепарабелен: его множители сепарабельны (\rf{2.24}), а различные
унитарные неприводимые многочлены не имеют общих корней (\rf{2.17}).

Пусть $f$ сепарабелен степени $n$, и $f(X)=c\prod_{i=1}^n(X-\alpha_i)$ в $F_f[X]$, $\alpha_i$ попарно различны. Каждый
$\sigma\in G_f$ переводит корень в корень: $f(\sigma\alpha_i)=\sigma(f(\alpha_i))=0$, так как коэффициенты $f$ лежат в $F$.
Будучи инъективным, $\sigma$ переставляет множество корней $\{\alpha_1,\dots,\alpha_n\}$. Так как $F_f=F[\alpha_1,\dots,\alpha_n]$,
автоморфизм $\sigma$ однозначно определяется этой перестановкой, а $(\sigma\tau)(\alpha_i)=\sigma(\tau(\alpha_i))$. Значит,
$\sigma\mapsto\sigma|_{\{\alpha_1,\dots,\alpha_n\}}$~--- инъективный гомоморфизм
\[
G_f\hookrightarrow\Sym(\{\alpha_1,\dots,\alpha_n\})\cong S_n ,
\]
и $G_f$ можно считать подгруппой симметрической группы. Изоморфизм с $S_n$ зависит от нумерации корней; другая нумерация
заменяет образ $G_f$ сопряжённой подгруппой $S_n$. Какие именно перестановки корней лежат в $G_f$, описывает следующее
предложение~--- это описание, не использующее ни полей, ни абстрактных групп (ни тех, ни других у Галуа не было).

\begin{proposition}[3.27b]
Перестановка $\sigma$ корней $\alpha_1,\dots,\alpha_n$ сепарабельного многочлена $f$ принадлежит $G_f$ тогда и только тогда,
когда она сохраняет все полиномиальные соотношения между корнями: для всех $P\in F[X_1,\dots,X_n]$
\begin{equation}\label{eq:c3-rel8}
P(\alpha_1,\dots,\alpha_n)=0\ \Longrightarrow\ P(\sigma\alpha_1,\dots,\sigma\alpha_n)=0 .
\end{equation}
\end{proposition}

\begin{proof}
Если $\tau\in G_f$, то $P(\tau\alpha_1,\dots,\tau\alpha_n)=\tau\bigl(P(\alpha_1,\dots,\alpha_n)\bigr)$, так как $\tau$~---
гомоморфизм колец, оставляющий на месте коэффициенты $P$; поэтому из $P(\alpha_1,\dots,\alpha_n)=0$ следует
$P(\tau\alpha_1,\dots,\tau\alpha_n)=\tau(0)=0$.

Обратно, пусть перестановка $\sigma$ удовлетворяет \eqref{eq:c3-rel8}. Рассмотрим гомоморфизм подстановки
\begin{equation}\label{eq:c3-rel9}
\varphi\colon F[X_1,\dots,X_n]\to F_f,\qquad X_i\mapsto\alpha_i .
\end{equation}
Его образ~--- кольцо $F[\alpha_1,\dots,\alpha_n]$, т.\,е.\ всё поле $F_f$ (подкольцо конечного расширения $F_f$, содержащее $F$,
является полем, \rf{1.23}), а ядро $I$~--- это в точности множество соотношений $P$ с
$P(\alpha_1,\dots,\alpha_n)=0$. Определим $s\colon F_f\to F_f$ формулой
\[
s\bigl(P(\alpha_1,\dots,\alpha_n)\bigr)=P(\sigma\alpha_1,\dots,\sigma\alpha_n),\qquad P\in F[X_1,\dots,X_n].
\]
\emph{Корректность}: каждый элемент $F_f$ имеет вид $P(\alpha_1,\dots,\alpha_n)$, но запись не единственна; если
$P(\alpha)=Q(\alpha)$, то $P-Q\in I$, по \eqref{eq:c3-rel8} $(P-Q)(\sigma\alpha)=0$, т.\,е.\ $P(\sigma\alpha)=Q(\sigma\alpha)$.
Именно здесь используется условие \eqref{eq:c3-rel8}: оно означает, что ядро $I$ гомоморфизма \eqref{eq:c3-rel9}
содержится в ядре гомоморфизма $X_i\mapsto\sigma\alpha_i$, и потому последний пропускается через
$F[X_1,\dots,X_n]/I\cong F_f$:
\[
\begin{tikzcd}[column sep=large]
F[X_1,\dots,X_n]\arrow[r,"X_i\mapsto\alpha_i",cold]\arrow[rd,"X_i\mapsto\sigma\alpha_i"',cold] & F_f\arrow[d,dashed,cnew,"s"]\\
 & F_f
\end{tikzcd}
\]
Так определённое $s$~--- гомоморфизм колец (значения сумм и произведений многочленов~--- суммы и произведения
значений), тождественный на $F$, и $s(\alpha_i)=\sigma\alpha_i$. Гомоморфизм полей автоматически инъективен (его
ядро~--- идеал поля, не содержащий $1$) и сохраняет деление: из $s(a/b)\,s(b)=s(a)$ следует $s(a/b)=s(a)/s(b)$. Наконец,
$s$ сюръективен: его образ~--- подполе, содержащее $F$ и все $\sigma\alpha_i$, т.\,е.\ все корни $\alpha_j$ (так как $\sigma$~---
перестановка), а они порождают $F_f$. Итак, $s\in\Aut(F_f/F)=G_f$, и $s$ действует на корнях как $\sigma$.
\end{proof}

Следствия.
\begin{enumerate}[label=(\alph*)]
\item Перестановки со свойством \eqref{eq:c3-rel8} автоматически сохраняют соотношения и в обратную сторону: $\sigma^{-1}$
тоже лежит в $G_f$ и потому удовлетворяет \eqref{eq:c3-rel8}. Так что $G_f$~--- это в точности перестановки, при
которых множество соотношений между корнями переходит в себя. (Проверять \eqref{eq:c3-rel8} достаточно для
каких-нибудь образующих идеала $I$: если $P=\sum Q_kP_k$ и все $P_k(\sigma\alpha)=0$, то и $P(\sigma\alpha)=0$.)
\item $|G_f|=[F_f:F]$ (\rf{3.17}), и по теореме Лагранжа $[F_f:F]$ делит $n!$~--- это усиливает оценку
$[F_f:F]\le n!$ из \rf{2.7}.
\end{enumerate}

\begin{example}[3.27c]
Пусть $f=X^4+1\in\Q[X]$, $\zeta=e^{\pi i/4}$, и корни занумерованы так: $\alpha_1=\zeta$, $\alpha_2=\zeta^3$, $\alpha_3=\zeta^5$,
$\alpha_4=\zeta^7$. Между ними есть соотношения
\[
\alpha_1+\alpha_3=0,\qquad\alpha_2+\alpha_4=0,\qquad\alpha_1\alpha_2+1=0\quad(\zeta^4=-1).
\]
Транспозиция $(\alpha_1\,\alpha_3)$ не лежит в $G_f$: она переводит третье соотношение в $\alpha_3\alpha_2+1=\zeta^8+1=2\ne0$.
Цикл $(\alpha_1\,\alpha_2\,\alpha_3\,\alpha_4)$ тоже не лежит в $G_f$: $\alpha_2\alpha_3+1=2\ne0$. Сама группа $G_f$ состоит из
автоморфизмов $\sigma_b\colon\zeta\mapsto\zeta^b$, $b\in\{1,3,5,7\}$ (\rf{5.8}, \rf{5.11g}), т.\,е.\ из перестановок
\[
\id,\quad(\alpha_1\,\alpha_2)(\alpha_3\,\alpha_4),\quad(\alpha_1\,\alpha_3)(\alpha_2\,\alpha_4),\quad(\alpha_1\,\alpha_4)(\alpha_2\,\alpha_3)
\]
(например, $\sigma_3$ переводит показатели $1,3,5,7$ в $3,9\equiv1,15\equiv7,21\equiv5$). Это четверная группа Клейна
порядка $4=[\Q[\zeta]:\Q]$, а не вся $S_4$.
\end{example}

\begin{proposition}[3.27d]
Пусть $f\in F[X]$ сепарабелен. Орбиты группы $G_f$ на множестве корней $f$~--- это в точности множества корней
неприводимых множителей $f$. В частности, $G_f$ действует на корнях транзитивно тогда и только тогда, когда $f$
неприводим, и тогда $\deg f$ делит $|G_f|$ \textup(ср.\ \rf{4.5}\textup).
\end{proposition}

\begin{proof}
Пусть $g$~--- неприводимый множитель $f$. Каждый $\sigma\in G_f$ переводит корни $g$ в корни $g$ (как выше: коэффициенты
$g$ лежат в $F$). Обратно, пусть $\alpha,\beta$~--- корни $g$. Существует $F$-изоморфизм $\varphi_0\colon F[\alpha]\to F[\beta]$,
$\alpha\mapsto\beta$ (\rf{2.1}: оба поля изоморфны $F[X]/(g)$). Поле $F_f$ является полем разложения $f$ и над $F[\alpha]$, и
над $F[\beta]$ (оно порождено корнями $f$), а $\varphi_0f=f$. По \rf{2.13a} $\varphi_0$ продолжается до изоморфизма
$\sigma\colon F_f\to F_f$, т.\,е.\ до элемента $\sigma\in G_f$ с $\sigma\alpha=\beta$. Если $f$ неприводим степени $n$, то
орбита корня $\alpha$ состоит из $n$ элементов, и $|G_f|=n\cdot|\Stab(\alpha)|$, где $\Stab(\alpha)=\Gal(F_f/F[\alpha])$.
\end{proof}

\subsection*{Группы Галуа приводимых многочленов}

Пусть теперь $f$ сепарабелен, но приводим. Как связана $G_f$ с группами Галуа множителей? Например, верно ли, что
группа Галуа произведения есть произведение групп Галуа? Мы увидим, что $G_f$ лишь \emph{вкладывается} в произведение,
причём каждый элемент группы Галуа множителя продолжается до элемента $G_f$, но не все наборы таких элементов
продолжаются одновременно.

Пусть $f=f_1f_2\cdots f_k$, где многочлены $f_i\in F[X]$ сепарабельны и попарно взаимно просты (например, $f_i$~---
различные неприводимые множители $f$; но можно брать и $f=gh$ с любыми взаимно простыми сепарабельными $g,h$).
Обозначим через $R_i$ множество корней $f_i$ в $F_f$, через $E_i=F[R_i]\subset F_f$~--- поле, порождённое ими; это поле
разложения $f_i$ над $F$, т.\,е.\ расширение Галуа поля $F$ с группой $G_{f_i}$, и $F_f=E_1E_2\cdots E_k$ (композит:
$F_f$ порождено всеми корнями $f$, т.\,е.\ объединением множеств $R_i$).

\begin{proposition}[3.27e]
В этих обозначениях:
\begin{enumerate}[label=\textup{(\alph*)}]
\item отображение $r\colon G_f\to G_{f_1}\times\dots\times G_{f_k}$, $\sigma\mapsto(\sigma|_{E_1},\dots,\sigma|_{E_k})$,~---
инъективный гомоморфизм, и каждая его компонента $G_f\to G_{f_i}$ сюръективна; на языке перестановок $G_f$ лежит в
подгруппе $\Sym(R_1)\times\dots\times\Sym(R_k)$ группы $\Sym(R_1\cup\dots\cup R_k)$;
\item $r$~--- изоморфизм тогда и только тогда, когда $[F_f:F]=\prod_i[E_i:F]$, что равносильно условию
\[
E_i\cap(E_1\cdots E_{i-1})=F\qquad\text{для }i=2,\dots,k;
\]
\item при $k=2$ образ $r$ равен
\[
\{(\sigma_1,\sigma_2)\in G_{f_1}\times G_{f_2}\mid \sigma_1|_{E_1\cap E_2}=\sigma_2|_{E_1\cap E_2}\}
=G_{f_1}\times_{\Gal(E_1\cap E_2/F)}G_{f_2},
\]
т.\,е.\ $G_{f_1f_2}$~--- расслоённое произведение групп $G_{f_1}$ и $G_{f_2}$ над $\Gal(E_1\cap E_2/F)$; в частности,
$G_{f_1f_2}\cong G_{f_1}\times G_{f_2}$ тогда и только тогда, когда $E_1\cap E_2=F$ \textup(\rf{3.21a}\textup).
\end{enumerate}
\end{proposition}

\begin{proof}
(a) Каждый $\sigma\in G_f$ переставляет корни каждого $f_i$, поэтому $\sigma(E_i)=E_i$, и $\sigma|_{E_i}\in\Aut(E_i/F)=G_{f_i}$.
Ограничение перестановочно с композицией, так что $r$~--- гомоморфизм. Если все $\sigma|_{E_i}$ тождественны, то $\sigma$
оставляет на месте все корни $f$, которые порождают $F_f$, и $\sigma=\id$. Сюръективность компоненты: поле $F_f$~---
поле разложения $f$ над $E_i$, и автоморфизм $\sigma_i\in G_{f_i}$ поля $E_i$ переводит $f$ в себя (коэффициенты в $F$);
по \rf{2.13a} он продолжается до автоморфизма $\sigma$ поля $F_f$, т.\,е.\ $\sigma\in G_f$ и $\sigma|_{E_i}=\sigma_i$. (Это же
следует из основной теоремы: $G_f\to\Gal(E_i/F)$~--- факторотображение по $\Gal(F_f/E_i)$, \rf{3.17}.)

(b) Так как $r$ инъективен, он является изоморфизмом тогда и только тогда, когда $|G_f|=\prod_i|G_{f_i}|$, т.\,е.\
$[F_f:F]=\prod_i[E_i:F]$. Положим $C_i=E_1\cdots E_i$. Поскольку $E_i/F$~--- расширение Галуа, по \rf{3.20}
\[
[C_i:F]=\frac{[C_{i-1}:F]\,[E_i:F]}{[C_{i-1}\cap E_i:F]},\qquad\text{откуда}\qquad
[F_f:F]=[C_k:F]=\frac{\prod_{i=1}^k[E_i:F]}{\prod_{i=2}^k[C_{i-1}\cap E_i:F]} .
\]
Равенство $[F_f:F]=\prod_i[E_i:F]$ выполняется тогда и только тогда, когда все $[C_{i-1}\cap E_i:F]=1$.

(c) Это предложение \rf{3.21}, применённое к расширениям Галуа $E_1,E_2$ с композитом $E_1E_2=F_f$.
\end{proof}

\begin{center}
\begin{keepfig}
\centering
\begin{tikzpicture}[every node/.style={font=\small},baseline=(current bounding box.center)]
\node (F) at (0,0) {$F$}; \node (I) at (0,1.1) {$E_1\cap E_2$};
\node (E1) at (-1.6,2.2) {$E_1=F_{f_1}$}; \node (E2) at (1.6,2.2) {$E_2=F_{f_2}$}; \node (C) at (0,3.3) {$F_f=E_1E_2$};
\draw (F) -- (I) -- (E1) -- (C) -- (E2) -- (I);
\end{tikzpicture}\qquad\qquad
\begin{tikzcd}[column sep=small,row sep=large,baseline=(current bounding box.center)]
 & G_f\arrow[ld,cnew,"\sigma\mapsto\sigma|_{E_1}"']\arrow[rd,cnew,"\sigma\mapsto\sigma|_{E_2}"] & \\
G_{f_1}\arrow[rd,cold] & & G_{f_2}\arrow[ld,cold]\\
 & \Gal(E_1\cap E_2/F) &
\end{tikzcd}
\figcap{Предложение \rf{3.27e}(c). Слева: поля разложения множителей и их пересечение. Справа: $G_f$ вкладывается в
$G_{f_1}\times G_{f_2}$, и её образ~--- пары $(\sigma_1,\sigma_2)$ с одинаковым образом в $\Gal(E_1\cap E_2/F)$ (ограничения на
$E_1\cap E_2$), т.\,е.\ расслоённое произведение.}
\end{keepfig}
\end{center}

\begin{remark}[3.27f: почему не всякий набор поднимается]
Попробуем по набору $(\sigma_1,\dots,\sigma_k)\in\prod G_{f_i}$ определить перестановку $\sigma$ всех корней <<покомпонентно>>:
$\sigma|_{R_i}=\sigma_i|_{R_i}$. По \rf{3.27b} $\sigma\in G_f$ тогда и только тогда, когда $\sigma$ сохраняет все соотношения
между корнями $f$. Соотношения между корнями одного множителя $f_i$ сохраняются (\rf{3.27b} для $f_i$, так как
$\sigma_i\in G_{f_i}$). Но бывают и <<смешанные>> соотношения, связывающие корни разных множителей, и именно они
препятствуют поднятию. Они возникают из общих элементов полей: если $\gamma\in E_1\cap E_2$, то
$\gamma=P_1(\text{корни }f_1)=P_2(\text{корни }f_2)$ для некоторых многочленов $P_1,P_2$ с коэффициентами из $F$, и
соотношение $P_1-P_2=0$ сохраняется тогда и только тогда, когда $\sigma_1\gamma=\sigma_2\gamma$. Это и есть условие
согласованности в \rf{3.27e}(c).
\end{remark}

\begin{example}[3.27g]
(a) $f=(X^2-2)(X^2-3)$ над $\Q$. Здесь $E_1\cap E_2=\Q[\sqrt2]\cap\Q[\sqrt3]=\Q$: иначе $\sqrt3=a+b\sqrt2$ с $a,b\in\Q$, и
$3=a^2+2b^2+2ab\sqrt2$ даёт $ab=0$, т.\,е.\ $3=a^2$ или $3=2b^2$~--- невозможно в $\Q$. Поэтому
$G_f\cong\Z/2\times\Z/2$; при нумерации $\alpha_{1,2}=\pm\sqrt2$, $\alpha_{3,4}=\pm\sqrt3$ это подгруппа
$\{\id,(\alpha_1\,\alpha_2),(\alpha_3\,\alpha_4),(\alpha_1\,\alpha_2)(\alpha_3\,\alpha_4)\}\subset S_4$ с двумя орбитами, как и
положено по \rf{3.27d}.

(b) $f=(X^3-2)(X^2+3)$ над $\Q$ (ср.\ пример после \rf{3.21a}). Пусть $\omega=e^{2\pi i/3}=\frac{-1+\sqrt{-3}}2$; корни $X^3-2$~--- это $a_1=\sqrt[3]2$,
$a_2=\omega\sqrt[3]2$, $a_3=\omega^2\sqrt[3]2$, и $E_1=\Q[\sqrt[3]2,\omega]$ имеет степень $6$ (вещественное поле $\Q[\sqrt[3]2]$ степени $3$ не
содержит $\omega$). Поэтому $G_{f_1}$, вложенная в $S_3$, совпадает с $S_3$. Корни $X^2+3$~--- это $b_1=\sqrt{-3}=\omega-\omega^2$ и $b_2=-b_1$, так что
$E_2=\Q[\sqrt{-3}]\subset E_1$. Значит, $E_1\cap E_2=E_2$, $F_f=E_1$ и $G_f\cong S_3$ имеет порядок $6$, а не $12=|S_3\times\Z/2|$.
Образ $r$ состоит из пар $(\sigma,\varepsilon(\sigma))$, где $\varepsilon(\sigma)=\sigma|_{E_2}$: $3$-цикл
$a_1\mapsto a_2\mapsto a_3\mapsto a_1$ переводит $\omega=a_2/a_1$ в $a_3/a_2=\omega$, а транспозиция $(a_2\,a_3)$ переводит
$\omega$ в $a_3/a_1=\omega^2$. Гомоморфизмы $\varepsilon$ и $\sgn$ из $S_3$ в $\{\pm1\}$ совпадают на этих двух
элементах, а они порождают $S_3$; поэтому $\varepsilon=\sgn$: $\sigma$ меняет знак $\sqrt{-3}$ ровно для нечётных $\sigma$.
Смешанное соотношение, отвечающее общему элементу $\sqrt{-3}$:
\[
a_1b_1=\sqrt[3]2\,(\omega-\omega^2)=a_2-a_3 .
\]
Покомпонентная перестановка <<$a_i$ на месте, $b_1\leftrightarrow b_2$>> его нарушает: $a_1b_2=-(a_2-a_3)\ne a_2-a_3$. А
перестановка $(a_2\,a_3)(b_1\,b_2)$ его сохраняет: $a_1b_2=a_3-a_2$.

(c) $f=(X^2-2)(X^2-3)(X^2-6)$ над $\Q$. Попарные пересечения полей $\Q[\sqrt2]$, $\Q[\sqrt3]$, $\Q[\sqrt6]$ равны $\Q$ (как
в (a)), но $\sqrt6=\sqrt2\sqrt3\in E_1E_2$, так что $E_3\cap E_1E_2=E_3\ne\Q$, и $G_f\cong(\Z/2)^2$ имеет порядок $4$, а не
$8$. Образ $r$ в $(\Z/2)^3=\{\pm1\}^3$ (знаки, на которые умножаются $\sqrt2,\sqrt3,\sqrt6$)~--- это тройки
$(\varepsilon_1,\varepsilon_2,\varepsilon_1\varepsilon_2)$; смешанное соотношение~--- $\sqrt2\cdot\sqrt3=\sqrt6$. Так что
\emph{попарной} тривиальности пересечений в \rf{3.27e}(b) недостаточно.

(d) $f=(X^3-2)(X^3-3)$ над $\Q$. Поля $E_1=\Q[\sqrt[3]2,\omega]$ и $E_2=\Q[\sqrt[3]3,\omega]$ степени $6$ содержат $\Q[\omega]$.
Пересечение $E_1\cap E_2$~--- расширение Галуа поля $\Q$ (\rf{3.21}), т.\,е.\ неподвижное поле нормальной подгруппы
$\Gal(E_1/\Q)\cong S_3$: одно из полей $\Q$, $\Q[\omega]$, $E_1$. Покажем, что $E_1\ne E_2$; тогда $E_1\cap E_2=\Q[\omega]$. Если
$E_2=E_1$, то $\Q[\sqrt[3]3]$~--- кубическое подполе $E_1$, т.\,е.\ неподвижное поле подгруппы порядка $2$; таких подполей
три: $\Q[a_1],\Q[a_2],\Q[a_3]$, и вещественно из них только $\Q[a_1]=\Q[\sqrt[3]2]$. Значит, $\sqrt[3]3=x+y\sqrt[3]2+z\sqrt[3]4$
с $x,y,z\in\Q$. Вложение $\tau\colon\Q[\sqrt[3]2]\to\CC$, $\sqrt[3]2\mapsto\omega\sqrt[3]2$, переводит $\sqrt[3]3$ в кубический
корень из $3$, т.\,е.\ в $\omega^l\sqrt[3]3$:
\[
x+y\omega\sqrt[3]2+z\omega^2\sqrt[3]4=\omega^l\bigl(x+y\sqrt[3]2+z\sqrt[3]4\bigr).
\]
Числа $1,\sqrt[3]2,\sqrt[3]4$ линейно независимы над $\Q[\omega]$ (так как $[E_1:\Q[\omega]]=3$), и сравнение
коэффициентов даёт $x=\omega^lx$, $y\omega=\omega^ly$, $z\omega^2=\omega^lz$; значит, ненулевым может быть лишь одно из
чисел $x,y,z$, и $3$ равно $x^3$, $2y^3$ или $4z^3$, что невозможно для рациональных $x,y,z$. Итак, $E_1\cap E_2=\Q[\omega]$,
$[F_f:\Q]=6\cdot6/2=18$ (\rf{3.20}), и
\[
G_f\cong S_3\times_{\Z/2}S_3=\{(\sigma,\tau)\in S_3\times S_3\mid\sgn\sigma=\sgn\tau\}
\]
(ограничение на $\Q[\omega]=\Q[\sqrt{-3}]$ задаётся знаком перестановки, как в (b)).
\end{example}

\begin{corollary}[3.27h]
В обозначениях \rf{3.27e} группа $G_f$ разрешима тогда и только тогда, когда разрешимы все группы $G_{f_i}$.
\end{corollary}

\begin{proof}
Если все $G_{f_i}$ разрешимы, то разрешимо их произведение, а значит, и его подгруппа $r(G_f)\cong G_f$. Обратно, каждая
$G_{f_i}$~--- гомоморфный образ $G_f$ (\rf{3.27e}(a)), а факторгруппы разрешимых групп разрешимы.
\end{proof}
